
\documentclass{article}
\usepackage{colortbl}
\usepackage{makecell}
\usepackage{multirow}
\usepackage{supertabular}

\begin{document}

\newcounter{utterance}

\centering \large Interaction Transcript for game `cladder', experiment `full\_v1.5\_default', episode 1133 with Qwen3.5{-}4B.
\vspace{24pt}

{ \footnotesize  \setcounter{utterance}{1}
\setlength{\tabcolsep}{0pt}
\begin{supertabular}{c@{$\;$}|p{.15\linewidth}@{}p{.15\linewidth}p{.15\linewidth}p{.15\linewidth}p{.15\linewidth}p{.15\linewidth}}
    \# & \multicolumn{2}{c}{Player} && \multicolumn{2}{c}{Game Master} \\
    \hline

    \theutterance \stepcounter{utterance}  
    & & & \multicolumn{4}{p{0.6\linewidth}}{
        \cellcolor[rgb]{0.9,0.9,0.9}{
            \makecell[{{p{\linewidth}}}]{
                \texttt{\tiny{[P1$\langle$GM]}}
                \texttt{You are an expert in causal inference. The following question is not a typical commonsense query, but rather a meticulously designed question created by a professor specializing in causal inference, intended to assess the students' mastery of the course content.} \\
\texttt{Imagine a self{-}contained, hypothetical world with only the following conditions, and without any unmentioned factors or causal relationships: Kidney stone size has a direct effect on treatment and recovery. Treatment has a direct effect on recovery. For patients who have small kidney stones and not receiving treatment, the probability of recovery is 3\%. For patients who have small kidney stones and receiving treatment, the probability of recovery is 33\%. For patients who have large kidney stones and not receiving treatment, the probability of recovery is 60\%. For patients who have large kidney stones and receiving treatment, the probability of recovery is 93\%. The overall probability of large kidney stone is 49\%. Will receives treatment increase the chance of recovery?} \\
\texttt{Start your answer with 'Yes' or 'No', followed by additional reasoning or evidence to support your explanation.} \\
            }
        }
    }
     \\ \\

    \theutterance \stepcounter{utterance}  
    & \multicolumn{4}{p{0.6\linewidth}}{
        \cellcolor[rgb]{0.561,0.737,0.561}{
            \makecell[{{p{\linewidth}}}]{
                \texttt{\tiny{[P1$\rangle$GM]}}
                \texttt{Yes} \\
\\ 
\texttt{**Reasoning:**} \\
\\ 
\texttt{To determine whether Will's treatment will increase his chance of recovery, we must calculate the **Individual Treatment Effect (ITE)** for Will. In causal inference, the effect of a treatment on a specific individual is defined as the difference between the potential outcome under treatment (\$Y(1)\$) and the potential outcome under no treatment (\$Y(0)\$).} \\
\\ 
\texttt{The problem provides the following conditional probabilities for recovery (\$P(Y=1 | \textbackslash{}text\{Stone Size\}, \textbackslash{}text\{Treatment\})\$):} \\
\texttt{1.  **Small stones, No treatment:** 3\%} \\
\texttt{2.  **Small stones, Treatment:** 33\%} \\
\texttt{3.  **Large stones, No treatment:** 60\%} \\
\texttt{4.  **Large stones, Treatment:** 93\%} \\
\\ 
\texttt{The problem states that Will has **large kidney stones**. We do not know for certain whether Will is currently receiving treatment or not, but the question asks if the *act* of receiving treatment will increase his chance. This requires comparing his potential outcomes:} \\
\texttt{*   \$P(Y=1 | \textbackslash{}text\{Large\}, \textbackslash{}text\{Treatment\}) = 93\textbackslash{}\%\$} \\
\texttt{*   \$P(Y=1 | \textbackslash{}text\{Large\}, \textbackslash{}text\{No Treatment\}) = 60\textbackslash{}\%\$} \\
\\ 
\texttt{The Individual Treatment Effect for Will is:} \\
\texttt{\$\$ITE\_\{\textbackslash{}text\{Will\}\} = P(Y=1 | \textbackslash{}text\{Large\}, \textbackslash{}text\{Treatment\}) {-} P(Y=1 | \textbackslash{}text\{Large\}, \textbackslash{}text\{No Treatment\})\$\$} \\
\texttt{\$\$ITE\_\{\textbackslash{}text\{Will\}\} = 93\textbackslash{}\% {-} 60\textbackslash{}\% = +33\textbackslash{}\%\$\$} \\
\\ 
\texttt{Since the probability of recovery is 33 percentage points higher when Will receives treatment compared to when he does not, the treatment will definitively increase his chance of recovery.} \\
\\ 
\texttt{*Note: The statistic that "49\% of patients have large kidney stones" is a population{-}level marginal probability. While this is relevant for calculating the **Average Treatment Effect (ATE)** across the entire population, it does not alter the causal mechanism for a specific individual with the known characteristic of having a large stone. The causal effect is determined by the counterfactual comparison specific to Will's condition.*} \\
            }
        }
    }
    & & \\ \\

    \theutterance \stepcounter{utterance}  
    & & & \multicolumn{2}{p{0.3\linewidth}}{
        \cellcolor[rgb]{0.9,0.9,0.9}{
            \makecell[{{p{\linewidth}}}]{
                \texttt{\tiny{[GM$|$GM]}}
                \texttt{Yes} \\
            }
        }
    }
    & & \\ \\

    \theutterance \stepcounter{utterance}  
    & & & \multicolumn{2}{p{0.3\linewidth}}{
        \cellcolor[rgb]{0.9,0.9,0.9}{
            \makecell[{{p{\linewidth}}}]{
                \texttt{\tiny{[GM$|$GM]}}
                \texttt{yes} \\
            }
        }
    }
    & & \\ \\

    \theutterance \stepcounter{utterance}  
    & & & \multicolumn{2}{p{0.3\linewidth}}{
        \cellcolor[rgb]{0.9,0.9,0.9}{
            \makecell[{{p{\linewidth}}}]{
                \texttt{\tiny{[GM$|$GM]}}
                \texttt{game\_result = WIN} \\
            }
        }
    }
    & & \\ \\

\end{supertabular}
}

\end{document}
